\section{Introduction}
\label{sec:intro}

% 1.5-2 pages.

\subsection{Problem}
% The unit of work is a ticket; the desired output is a merged pull request.
% Producing the diff is increasingly tractable. Establishing that the diff is
% correct is not, and it is what actually gates deployment.

\subsection{Why existing orchestration does not fit}
% Two families, both mismatched:
%   - CI/CD and workflow engines (Argo, Tekton, Airflow, Temporal): assume the
%     tree is immutable for the duration of the run. Nodes execute once.
%   - Agent frameworks (LangGraph, OpenHands, SWE-agent): model the reasoning
%     loop but provide no isolation, no capacity model, no cluster primitives.
% Neither treats verification state as a first-class, re-entrant artifact, and
% neither separates the step that changes the tree from the step that judges it.

\subsection{Contributions}
% Enumerate 5. Draft:
%   (C1) A formulation of agent scaling as a verification-capacity problem, with
%        verification density as the tunable parameter.
%   (C2) A content-addressed verification ledger with default-closed,
%        path-scoped invalidation and a pure convergence core (\fold, \decide),
%        giving well-defined convergence and a computable termination bound for
%        a tree that mutates between gate runs.
%   (C3) A declarative workflow specification that separates capability
%        (cluster admin: what a step IS -- action or gate) from composition
%        (team: which steps, in what order, invalidated by what), with
%        integrity gates injected by policy at plan resolution so that no
%        workflow can compose them away or shadow them with a weaker copy.
%   (C4) A threat model for agent-authored code in a shared cluster, including
%        gate-gaming as adversarial selection pressure under best-of-N and the
%        observation that invalidation defaults are a security property.
%   (C5) An open-source Kubernetes implementation -- five CRDs, two controllers,
%        no admission webhooks -- and a measured cost datapoint.

\todo{Draft after Section 2 and Section 6 settle.}
